% --- CONCLUSION (sans numérotation) ---
\chapter*{Conclusion}
\addcontentsline{toc}{chapter}{Conclusion}
\label{chap:conclusion}

Ce travail est parti d'un objectif modeste : concevoir un assistant IA. Le choix du cas s'est porté sur SOYAFF, l'entreprise où je travaille comme informaticien. À mesure que le projet avançait, un problème plus profond est apparu. Les chefs de Bamako devaient décider sur des opérations menées à Dakar sans aucune visibilité en temps réel. Ils recevaient des fichiers Excel de plus d'une cinquantaine d'opérations, où l'information se noyait, et les incompréhensions se réglaient par de longs appels téléphoniques entre les deux bureaux. L'assistant est devenu une plateforme complète. Le chef logistique a donné son accord pour la construire, en posant ses conditions : la sécurité, la souveraineté des données. Presque toutes ont été traitées entre le développement et la mise en production.

\textbf{Ce que la plateforme a changé.} Une opération est désormais saisie une seule fois à Dakar et visible aussitôt à Bamako. Le temps de saisie a baissé d'environ 85~\%. Les utilisateurs ont attribué à la plateforme un score d'utilisabilité moyen de 88,1 sur 100. C'est le résultat dont je suis le plus fier : savoir que ce travail facilite le quotidien de mes collègues. Le tout fonctionne sur un seul serveur pour 37,4~€ par mois. Ce chiffre traduit un choix de conception : une PME ouest-africaine a un budget limité et n'est pas prête à payer lourdement pour sa digitalisation, quitte à rester sur Excel. L'architecture a été pensée pour cette réalité. En valorisant le temps libéré, l'investissement initial est amorti en moins de deux ans, et en moins d'un an si l'on compte la fin de la double saisie à Bamako.

\textbf{Ce que le benchmark m'a appris.} Le RAG local n'a répondu juste qu'à 20~\% des questions, contre 88,9~\% pour le cloud. Ce résultat a été une vraie déception. J'ai d'abord cru que le modèle Mistral~7B était en cause. En reprenant les réponses une à une, j'ai compris que je comparais deux méthodes de récupération différentes. L'agent cloud écrit une requête SQL sur toute la base ; le chemin local cherche des passages par similarité sémantique. Or cette recherche est, par construction, inadaptée à des données relationnelles et à des questions d'agrégat : aucun modèle ne peut calculer un total juste à partir de dix lignes approchantes. C'est la principale leçon d'ingénieur de ce travail. Elle vaut au-delà de SOYAFF : avant de comparer deux modèles, il faut s'assurer qu'ils reçoivent les données de la même façon.

\textbf{Ce qui reste à faire.} La suite logique est de doter le chemin local d'un agent \textit{text-to-SQL} exécuté sur le serveur, avec un modèle spécialisé dans la génération de requêtes, comme SQLCoder de Defog, Code Llama ou DeepSeek-Coder. Le chemin souverain disposerait alors de la méthode qui réussit côté cloud, et le benchmark pourrait être rejoué pour retester H2. La réplicabilité reste elle aussi à démontrer sur un second cas réel : l'analyse par lignes de produits logiciels a identifié ce qui se garde et ce qui se reconfigure, mais le déploiement chez SODRAF, entreprise partenaire de SOYAFF, n'a pas pu être mené faute de temps. Après la soutenance, une réunion à Bamako doit décider de la suite du projet avec la direction de SOYAFF.

La digitalisation des PME ouest-africaines reste un défi. Il revient à ses ingénieurs de le relever, pour garder la maîtrise de nos données et rendre nos entreprises plus productives. C'est notre façon de contribuer au développement de l'Afrique. Ce travail montre qu'on peut commencer avec un seul serveur, des outils libres et des utilisateurs qui disent franchement ce qui leur manque.
